\chapter{Introduction}\label{ch:introduction}

A production bottleneck is the resource or process that constrains the throughput of a manufacturing system at a particular time~\cite{subramaniyan2023throughput,tang2024comprehensive}. Bottleneck identification is consequently a system-level problem: a busy machine, a long local queue, or a temporarily unavailable resource does not necessarily determine the output of the entire production line. In a dynamic production environment, the throughput-constraining resource may also change as machines become unavailable, materials are delayed, work-in-process (WIP) accumulates, or operating conditions vary~\cite{lai2021data}. Production management therefore benefits not only from detecting the currently active constraint but also from predicting which resource may become constrained, when the event is likely to begin, and how long it may persist. Advance warning can provide information for resource allocation and production scheduling, although whether those decisions improve actual throughput must be evaluated separately.

Dynamic bottleneck prediction is especially challenging in a manufacturing system that contains multiple interdependent resource types. Workstations carry out processing and assembly; workers provide services; gantries and automated guided vehicles (AGVs) transport materials; and buffers temporarily store intermediate products. The effect of a local operating change can move through several of these entities before it becomes visible as a constraint elsewhere. For example, a downstream workstation may continue processing after an upstream interruption while its intermediate buffer still contains WIP. Only after that supply is exhausted might it become starved. Conversely, a downstream interruption may eventually block an upstream workstation once the intervening buffer fills. Studies of manufacturing buffers and disruption effects demonstrate that the origin of an interruption and the resource affected later need not coincide~\cite{gu2015passive,zhang2022vulnerability}. Historical evidence distributed across resources may thus be valuable before the eventual bottleneck displays an obvious local symptom.

Existing bottleneck-prediction research has progressed from time-series extrapolation of machine indicators to recurrent and graph-based models. Statistical methods and active-period analysis estimate future throughput constraints from operational histories~\cite{li2011throughput,subramaniyan2018data}, while recurrent networks learn temporal relationships in station-level observations~\cite{liu2023lstmdam}. Manufacturing graph methods additionally represent production relations among stations or tasks~\cite{lai2023bstan,yin2025btgcn}. Nevertheless, the joint representation of workstations, workers, material-supply conditions, buffers, and transport resources remains limited. A model observing primarily the target workstation may overlook earlier signals elsewhere in the production and logistics system. Although graph aggregation and temporal recurrence can represent useful dependencies, they do not by themselves ensure that an earlier cross-resource operating pattern is aligned with a later bottleneck event.

A second difficulty is the construction of reliable bottleneck supervision. Sustained bottleneck events are relatively scarce during normal production, and deliberately introducing failures, labour absences, shortages, or transport interruptions in a real factory is generally impractical. Operational indicators based on utilization, blockage, starvation, or queue behaviour support retrospective assessment, but isolated thresholds can confuse brief fluctuations with persistent constraints. Labels based on a single measurement can also obscure differences among observable operating conditions associated with an event~\cite{subramaniyan2023throughput,su2022dynamic}. A useful prediction target should therefore be defined through traceable process evidence and temporal persistence, without equating the source of a simulated disturbance with an established physical root cause.

This project addresses these limitations through \normalbold{PHGST-DM}, a process-informed heterogeneous-graph spatiotemporal transformer with delay-aware modelling. The framework uses traceable multi-entity production histories from an open-workspace NVIDIA Isaac Sim factory model. It represents workstations and handling resources within a heterogeneous process--logistics graph, while worker availability, buffer occupancy, material shortages, and transport conditions provide relevant operational context. Resource-specific rules and persistence requirements construct targets for sustained bottleneck events and observable process causes. Training-only temporal pattern clustering then summarizes recurrent operating behaviour, and a delay-aware spatiotemporal encoder uses this pattern context to condition cross-resource attention. The resulting predictions combine resource-level bottleneck occurrence, onset and duration with graph-level process-cause and remaining-order-time estimates.

The distinction between \emph{delay-aware prediction} and \emph{physical delay estimation} is important. PHGST-DM matches recent queue trajectories against activity-pattern prototypes and incorporates the matched context into geographic attention. This mechanism exploits potentially informative historical relationships for future-event prediction; it does not explicitly estimate the time required for a physical disturbance to propagate between two resources, nor do learned attention weights establish causality. Likewise, the process-cause categories describe observable operational symptoms and are not independently verified root-cause labels. These limits determine how the method and its results are interpreted throughout this report.

The project makes three methodological contributions. \normalbold{First}, it extends workstation-centred forecasting to a heterogeneous, multi-entity representation that combines production and logistics relationships with resource-specific operational states. \normalbold{Second}, it develops process-informed event and cause supervision, together with temporal pattern clustering, to distinguish sustained constraints from transient fluctuations and represent recurring operating behaviour. \normalbold{Third}, it introduces recurrent-pattern-guided delay-aware modelling that conditions cross-resource attention on matched historical patterns, with component evaluations examining its value for earlier bottleneck warning. The framework jointly reports the affected resource, onset, future duration, observable process cause, and remaining order-completion time. In the evaluated simulation setting, the final paper reports improvements in strict bottleneck-event F1 over the compared baselines; this does not imply that every auxiliary prediction objective is the best among all alternatives.

The remainder of the report is organized as follows. Chapter~\ref{ch:related} reviews bottleneck prediction, operational indicators, delayed production effects, and spatiotemporal graph learning. Chapter~\ref{ch:methodology} develops the problem formulation, data pipeline, supervision, graph representation, predictive architecture, training procedure, and production-sensitivity validation protocol. Chapter~\ref{ch:experiments} evaluates the method and its components. Chapter~\ref{ch:challenges} discusses the challenges encountered and how they were addressed, and Chapter~\ref{ch:conclusion} concludes the study and outlines future work.
